(1) 无穷远处 \(E=mv_\infty^2/2\)、\(L=mv_\infty b\)。近日点径向速度为零，故 \(E=L^2/(2mr_{\min}^2)+\kappa/r_{\min}\)；乘以 \(2mr_{\min}^2\) 后解二次方程，得 \(r_{\min}=[\kappa+\sqrt{\kappa^2+m^2v_\infty^4b^2}]/(mv_\infty^2)\)。只取正根，且 \(r_{\min}\to2\kappa/(mv_\infty^2)\) 为正面撞击的转折点。
(2) 对 \(u=1/r\) 使用 Binet 公式，排斥力为 \(F_r=\kappa u^2\)，从而 \(u''+u=-m\kappa/L^2\)。以近心点为极角原点，解为 \(u=(m\kappa/L^2)(e\cos\phi-1)\)，\(e=\sqrt{1+(mv_\infty^2 b/\kappa)^2}>1\)。无穷远处 \(u=0\) 的两条渐近方向满足 \(\cos\phi_\infty=1/e\)；轨道总偏转角为 \(\Theta=\pi-2\phi_\infty=2\arctan[\kappa/(mv_\infty^2b)]\)，处于 \(0<\Theta<\pi\)。
(3) 对轴对称中心势，冲量参数环的入射面积 \(d\sigma=2\pi b\,|db|\)，对应散射立体角 \(d\Omega=2\pi\sin\Theta\,d\Theta\)，因此 \(d\sigma/d\Omega=(b/\sin\Theta)|db/d\Theta|\)。由 \(b=[\kappa/(mv_\infty^2)]\cot(\Theta/2)\) 求导并整理，得到 \(d\sigma/d\Omega=[\kappa/(2mv_\infty^2)]^2\csc^4(\Theta/2)\)，即经典 Rutherford 型的角分布。
(4) \(b\to0^+\) 时 \(\Theta\to\pi\)，为强烈反向散射；\(b\to\infty\) 时 \(\Theta\sim2\kappa/(mv_\infty^2b)\to0\)，远处弱相互作用只引起小角偏转。对固定较大 \(\Theta\)，需要较小的冲量参数环面积，所以大角事件的测度比弱偏转小；若积分至零散射角，全截面因长程 \(1/r\) 势而发散，需物理屏蔽。
